\subsection{Схема 3: Восьмёркообразные узлы}

Восьмёрка образуется, если сделать колышку, но перед пробивкой обнести ходовой конец вокруг коренного
Восьмёрка получится, если сделать колышку, перекрутить её ещё на полоборота, а потом пробить.
Если накинуть оборот, получится стивидорный узел

Узлы, подобные восьмёрке содержат две группы пересечений из чётного числа каждое.

\graphicspath{{\currentpath}}

\includegraphics[scale=2]{images/stividor-0.eps}

\includegraphics[scale=2]{images/stividor-1.eps}

\includegraphics[scale=2]{images/stividor-2.eps}

\includegraphics[scale=2]{images/stividor-3.eps}

\subsection{Узлы похожие на восьмёрку, но с другой структурой}

Если сделать колышку и перекрутить её ещё на оборот, получится узел, похожий на стивидорный, но с нечётным числом пересечений.

Чтобы получить симметричный ему узел, нужно завязать простой узел, затем на сделать на нём колышку и пробить её ходовым концом.

Подобный узел может содержать только нечётное число пересечений
